\documentclass[main]{subfiles}


\title[Evaluation and Benchmarking]{Evaluation and Benchmarking}
\subtitle{Benchmarking of AutoML}
\author[Marcel Wever]{Marcel Wever}
\date{}

\titlebackground*{images/AutoML_ML_Evaluation}

%\institute{}
%\date{}

\begin{document}
	
	\maketitle
	

%----------------------------------------------------------------------
%----------------------------------------------------------------------
\begin{frame}[c]{Benchmarking of AutoML}
\begin{itemize}
    \item Benchmarking = systematic comparison of methods across a defined set of tasks
    \item Goal: draw more general conclusions about relative performance that go beyond a single dataset
    \item A benchmark consists of:
    \begin{itemize}
        \item A task suite: A collection of datasets or problem instances
        \item A protocol: fixed evaluation procedure, time budgets, and hardware specification
        \item A performance measure: the metric used to compare systems
    \end{itemize}
    \item Good benchmarks are reproducible, fair, and representative of real-world use
\end{itemize}
\end{frame}

\begin{frame}{No Algorithm Dominates All Others Across All Tasks}
    \begin{itemize}
        \item Consequence: the result of a comparison depends on the choice of task suite
        \item A benchmark implicitly defines a distribution over tasks, results only generalize to tasks drawn from a similar distribution
        \item \alert{Cherry-picking datasets} to favor one's own system is a serious threat to validity
        \item A credible benchmark should cover a diverse set of tasks:
        \begin{itemize}
            \item varying size
            \item feature types
            \item class imbalance
            \item difficulty
        \end{itemize}
    \end{itemize}
\end{frame}

\begin{frame}{Desiderata for a Fair Benchmark}
    \begin{itemize}
        \item Documentation: Important for adoption, correct use, explicit conventions
        \item Reproducibility: Fixed instances, public datasets, containerized environments
        \item Controlled resources: Same time budget, memory limit, and hardware for all systems
        \item Nested evaluation: Outer test folds held out, inner loop used for search
        \item Task diversity: Sufficient variety to avoid conclusions driven by a single dataset
        \item Baselines: Always include simple baselines (e.g., random search, single well-tuned model) to contextualize results
        \item Openness: Code, data splits, and raw results publicly available for scrutiny
    \end{itemize}
\end{frame}

\begin{frame}{Exiting AutoML Benchmarks}
    \begin{itemize}
        \item OpenML-CC18 \liturl{Bischl et al.\ 2021}{https://datasets-benchmarks-proceedings.neurips.cc/paper_files/paper/2021/file/c7e1249ffc03eb9ded908c236bd1996d-Paper-round2.pdf}: 72 curated classification datasets from OpenML, widely used as a default task suite
        \item AutoML Benchmark \liturl{Gijsbers et al.\ 2024}{https://www.jmlr.org/papers/volume25/22-0493/22-0493.pdf}: Standardized framework for comparing AutoML systems with controlled resources and reproducible splits
        \item HPOBench / HPO-B \liturl{Eggensperger et al. 2021}{https://datasets-benchmarks-proceedings.neurips.cc/paper_files/paper/2021/file/93db85ed909c13838ff95ccfa94cebd9-Paper-round2.pdf}: Focused benchmarks for hyperparameter optimization, providing surrogate evaluations or lookup tables for cheap repeated experiments
        \item YAHPO Gym \liturl{Pfisterer et al.\ 2022}{https://proceedings.mlr.press/v188/pfisterer22a/pfisterer22a.pdf}: Surrogate-based benchmark suite covering a large space of HPO scenarios efficiently
        \item PD1 \liturl{Wang et al.\ 2024}{https://jmlr.org/papers/volume25/23-0269/23-0269.pdf}: HPO for modern deep learning architectures (ResNets and Transformer)
        \item CARP-S \liturl{Benjamins et al.\ 2025}{https://arxiv.org/pdf/2506.06143}: Benchmarking framework   HPOBench, YAHPO Gym, PD1, and more
    \end{itemize}
\end{frame}

\begin{frame}{Aggregating Results Across Tasks}
    \begin{itemize}
        \item Raw performance scores are not directly comparable across datasets (different scales, metrics)
        \item Common normalization strategies:
        \begin{itemize}
            \item Rank-based: Rank systems per dataset, then average ranks, robust but loses magnitude information
            \item Normalized score: Scale performance relative to a baseline and a reference ceiling per dataset
        \end{itemize}
        \item After normalization, apply the statistical testing workflow from earlier:
        \begin{itemize}
            \item Friedman test for overall differences
            \item Post-hoc tests with correction (Nemenyi / Holm) for pairwise comparisons
            \item Visualize with critical difference diagrams
        \end{itemize}
    \end{itemize}
\end{frame}

\begin{frame}{Threats to Validity}
    \begin{itemize}
        \item Time budget sensitivity: Rankings can change substantially with different budgets, always report results at multiple budgets or as anytime curves
        \item Hardware heterogeneity: Wall-clock time is hardware-dependent; document and control carefully
        \item Dataset leakage: Public datasets may have been used during system development, giving an unfair advantage
        \item Metric choice: A system optimized for accuracy may rank differently when evaluated on AUC or log-loss
        \item Stochasticity: Single-run comparisons are insufficient; always use multiple runs and report uncertainty
    \end{itemize}
\end{frame}

\begin{frame}{Summary} 

Most important insights from this video:

\begin{enumerate}
    \item A benchmark is only as credible as its task suite, protocol, and openness
    \item Cherry-picking datasets undermines any conclusion
    \item Fair benchmarking requires controlled resources, nested evaluation, and multiple runs with reported uncertainty
    \item Raw scores across datasets must be normalized (e.g., via ranks) before aggregation, followed by appropriate statistical tests
    \item Benchmark results are always relative to a specific task distribution and time budget
    \item Generalization beyond that scope requires caution
\end{enumerate}

\end{frame}

\end{document}
